\documentclass{ximera}

\input{../../../preamble.tex}


\definecolor{fvdnavy}{HTML}{071D43}
\definecolor{fvdcyan}{HTML}{20BFC6}
\definecolor{fvdcyanlight}{HTML}{EFFCFB}
\definecolor{fvdmagenta}{HTML}{F10D69}
\definecolor{fvdmagentalight}{HTML}{FFF1F7}
\definecolor{fvdtext}{HTML}{163044}





%=================================================
% Pedagogische omgevingen
% PDF: tcolorbox
% HTML: learning-box
%=================================================

\newenvironment{voorkennis}
{%
    \ifdefined\HCode
        \HCode{%
<section class="learning-box learning-box--prior">
<div class="learning-box__header">
<span class="learning-box__icon" aria-hidden="true"></span>
<span class="learning-box__title">Voorkennis</span>
</div>
<div class="learning-box__content">
        }%
        \begin{itemize}
    \else
       \begin{tcolorbox}[
    enhanced,
    breakable,
    colback=fvdcyanlight,
    colframe=fvdcyan,
    coltext=fvdtext,
    boxrule=0.5pt,
    leftrule=4pt,
    arc=4mm,
    outer arc=4mm,
    left=7mm,
    right=6mm,
    top=5mm,
    bottom=5mm,
    before skip=6mm,
    after skip=6mm,
    title=Voorkennis,
    coltitle=fvdcyan,
    fonttitle=\bfseries\Large,
    attach title to upper={\par\smallskip},
    boxed title style={
        empty,
        left=0mm,
        right=0mm,
        top=0mm,
        bottom=0mm
    },
    overlay={
        \draw[
            fvdcyan,
            line width=0.5pt
        ]
        ([xshift=42mm,yshift=-13mm]frame.north west)
        --
        ([xshift=-7mm,yshift=-13mm]frame.north east);
    }
]
        \begin{itemize}
    \fi
}
{%
    \end{itemize}
    \ifdefined\HCode
        \HCode{%
</div>
</section>
        }%
    \else
        \end{tcolorbox}
    \fi
}


\newenvironment{leerdoelen}
{%
    \ifdefined\HCode
        \HCode{%
<section class="learning-box learning-box--goals">
<div class="learning-box__header">
<span class="learning-box__icon" aria-hidden="true"></span>
<span class="learning-box__title">Leerdoelen</span>
</div>
<div class="learning-box__content">
        }%
        \begin{itemize}
    \else
\begin{tcolorbox}[
    enhanced,
    breakable,
    colback=fvdmagentalight,
    colframe=fvdmagenta,
    coltext=fvdtext,
    boxrule=0.5pt,
    leftrule=4pt,
    arc=4mm,
    outer arc=4mm,
    left=7mm,
    right=6mm,
    top=5mm,
    bottom=5mm,
    before skip=6mm,
    after skip=6mm,
    title=Leerdoelen,
    coltitle=fvdmagenta,
    fonttitle=\bfseries\Large,
    attach title to upper={\par\smallskip},
    boxed title style={
        empty,
        left=0mm,
        right=0mm,
        top=0mm,
        bottom=0mm
    },
    overlay={
        \draw[
            fvdmagenta,
            line width=0.5pt
        ]
        ([xshift=42mm,yshift=-13mm]frame.north west)
        --
        ([xshift=-7mm,yshift=-13mm]frame.north east);
    }
]
        \begin{itemize}
    \fi
}
{%
    \end{itemize}
    \ifdefined\HCode
        \HCode{%
</div>
</section>
        }%
    \else
        \end{tcolorbox}
    \fi
}


\addPrintStyle{../../..}

\title{Oefeningen --- Samenhang is nog geen oorzaak}
\author{Ingmar Herreman}



\begin{document}

% FVD_CHAPTER_DATA_START
% Automatisch gegenereerd uit index.tex.
% Niet handmatig aanpassen.
\fvdchapterdata{3}{Verbanden onderzoeken met data}{11}
% FVD_CHAPTER_DATA_END






\maketitle




\begin{oefening}[Wat is werkelijk aangetoond?]{\oefB \ESS}
In een onderzoek bij $800$ leerlingen blijkt dat leerlingen die meer
uren per week sporten gemiddeld een hogere score behalen op een
conditietest.

\begin{question}
Welke uitspraak wordt rechtstreeks door deze informatie ondersteund?

\begin{multipleChoice}
  \choice{Sporten veroorzaakt bij iedere leerling een hogere conditie.}
  \choice[correct]{In de onderzochte groep hangen sporturen en score op de conditietest positief samen.}
  \choice{De conditietest veroorzaakt dat leerlingen meer sporten.}
  \choice{Er zijn geen andere verklaringen mogelijk.}
\end{multipleChoice}
\end{question}

\begin{question}
Waarom mag je uit de gegeven informatie nog geen definitieve causale
conclusie trekken?

\begin{oplossing}
We weten alleen dat beide grootheden samenhangen. Andere factoren
kunnen een rol spelen en we kennen de onderzoeksopzet nog niet.
\end{oplossing}
\end{question}
\end{oefening}


\begin{oefening}[De derde variabele]{\oefB \ESS}
Tijdens de zomer blijkt er een positieve samenhang te bestaan tussen
de dagelijkse ijsverkoop en het aantal bezoekers van een openluchtzwembad.

\begin{question}
Geef een plausibele derde variabele die beide grootheden kan beïnvloeden.

\begin{oplossing}
De buitentemperatuur. Op warmere dagen worden doorgaans meer ijsjes
verkocht en bezoeken meer mensen het zwembad.
\end{oplossing}
\end{question}

\begin{question}
Bewijst dit dat temperatuur werkelijk de volledige verklaring is?

\begin{oplossing}
Nee. Het is een plausibele alternatieve verklaring die toont dat de
oorspronkelijke samenhang geen bewijs is voor een directe causale
relatie.
\end{oplossing}
\end{question}
\end{oefening}


\begin{oefening}[Welke richting?]{\oefB \ESS}
Onderzoekers vinden een positieve samenhang tussen het aantal uren dat
mensen bewegen en hun lichamelijke conditie.

Geef twee mogelijke causale richtingen.

\begin{oplossing}
Bijvoorbeeld:
\[
\text{meer bewegen}\longrightarrow\text{betere conditie},
\]
maar ook
\[
\text{betere conditie}\longrightarrow\text{meer kunnen bewegen}.
\]
Beide richtingen kunnen bovendien tegelijk een rol spelen.
\end{oplossing}
\end{oefening}


\begin{oefening}[Correlatie is geen percentage oorzaak]{\oefB \ESS}
Een onderzoek geeft voor twee grootheden

\[
r=0{,}82.
\]

Een leerling schrijft:

\begin{quote}
``De ene grootheid veroorzaakt de andere voor $82\%$.''
\end{quote}

Beoordeel deze uitspraak.

\begin{oplossing}
De uitspraak is fout. $r=0{,}82$ beschrijft een sterke positieve
lineaire samenhang. De correlatiecoëfficiënt is geen percentage
causaliteit en toont op zichzelf geen oorzaak-gevolgrelatie aan.
\end{oplossing}
\end{oefening}





\begin{oefening}[Vaders en zonen: samenhang zonder eenvoudig oorzakelijk verhaal]{\oefU \ESS}

Gebruik

\csv
{https://raw.githubusercontent.com/Ingmar-1986/wiskunde_cursus_2/main/modules/Data/Statistiek/vaderzoon.csv}
{vaderzoon.csv}

Maak een spreidingsdiagram en bepaal met ICT $r$.

\begin{question}
Welke statistische conclusie wordt door de gegevens ondersteund?

\begin{oplossing}
Er is in deze dataset een positieve lineaire samenhang tussen de lengte
van vaders en die van hun zonen. De correlatie bedraagt ongeveer
\[
r=0{,}704.
\]
\end{oplossing}
\end{question}

\begin{question}
Waarom is de uitspraak
\begin{quote}
``Een lange vader veroorzaakt rechtstreeks een lange zoon''
\end{quote}
te eenvoudig?

\begin{oplossing}
Lichaamslengte hangt samen met erfelijke factoren én met tal van andere
biologische en omgevingsfactoren. De correlatie beschrijft het patroon
in de gegevens maar bewijst geen eenvoudige één-op-één oorzaak.
\end{oplossing}
\end{question}

\end{oefening}


\begin{oefening}[Examenscores en een verleidelijke conclusie]{\oefU \ESS}

Gebruik

\csv
{https://raw.githubusercontent.com/Ingmar-1986/wiskunde_cursus_2/main/modules/Data/Statistiek/examen.csv}
{examen.csv}

Maak een spreidingsdiagram en bepaal $r$.

\begin{question}
Beschrijf wat de gegevens over de twee examenscores aantonen.

\begin{oplossing}
De lineaire correlatie is zeer klein:
\[
r\approx-0{,}107.
\]
In deze dataset is dus nauwelijks lineaire samenhang zichtbaar tussen
de twee scores.
\end{oplossing}
\end{question}

\begin{question}
Mag je daaruit besluiten dat kennis voor het ene vak nooit iets met
het andere vak te maken heeft?

\begin{oplossing}
Nee. De dataset en $r$ geven alleen informatie over de waargenomen
lineaire samenhang tussen deze scores. Ze bewijzen geen algemene
uitspraak over kennis, leerprocessen of causaliteit.
\end{oplossing}
\end{question}

\end{oefening}


\begin{oefening}[Festivalgegevens: wat veroorzaakt de klachten?]{\oefU \ESS}

Gebruik

\csv
{https://raw.githubusercontent.com/Ingmar-1986/wiskunde_cursus_2/main/modules/Data/Statistiek/festival.csv}
{festival.csv}

De dataset vergelijkt afstand tot het festivalterrein en het aantal
klachten.

\begin{question}
Welke samenhang zie je globaal in de puntenwolk?

\begin{oplossing}
Op grotere afstand zijn er globaal minder klachten, al is de
puntenwolk onregelmatig en zijn er invloedrijke waarnemingen.
\end{oplossing}
\end{question}

\begin{question}
Noem minstens drie mogelijke factoren naast afstand die het aantal
klachten kunnen beïnvloeden.

\begin{oplossing}
Bijvoorbeeld geluidsniveau en windrichting, bebouwing en afscherming,
bevolkingsdichtheid, tijdstip, gevoeligheid van bewoners of de ligging
ten opzichte van luidsprekers. Andere goed gemotiveerde antwoorden zijn
mogelijk.
\end{oplossing}
\end{question}

\begin{question}
Formuleer een causale conclusie die \emph{niet} door deze dataset alleen
gerechtvaardigd is.

\begin{oplossing}
Bijvoorbeeld: ``Elke extra kilometer afstand veroorzaakt een daling
van het aantal klachten.'' De gegevens tonen een samenhang, maar de
onderzoeksopzet sluit andere verklaringen niet uit.
\end{oplossing}
\end{question}

\end{oefening}


\begin{oefening}[Een betere proef]{\oefU \ESS}

Een leerling wil onderzoeken of een nieuw type lamp de groei van
bonenplanten beïnvloedt.

Hij zet tien planten onder de nieuwe lamp in een warme kamer en tien
andere planten onder een gewone lamp in een koelere kamer.

Na drie weken zijn de planten onder de nieuwe lamp gemiddeld groter.

\begin{question}
Waarom kan hij nog niet besluiten dat het nieuwe type lamp de grotere
groei veroorzaakte?

\begin{oplossing}
Naast het type lamp verschilde ook de temperatuur. Daardoor weten we
niet welke factor het verschil in groei verklaart.
\end{oplossing}
\end{question}

\begin{question}
Hoe zou je het experiment verbeteren?

\begin{oplossing}
Gebruik vergelijkbare planten, verdeel ze bij voorkeur willekeurig over
de groepen en houd relevante omstandigheden zoals temperatuur, water,
bodem, potgrootte en belichtingstijd gelijk. Laat alleen het onderzochte
lamptype systematisch verschillen.
\end{oplossing}
\end{question}

\end{oefening}




\begin{oefening}[Koppen kritisch lezen]{\oefU \ESS}

Een nieuwsbericht heeft als titel:

\begin{quote}
``Ontbijt maakt leerlingen slimmer.''
\end{quote}

In het beschreven onderzoek werden $1200$ leerlingen bevraagd.
Leerlingen die regelmatig ontbijten, behalen gemiddeld hogere
schoolresultaten.

\begin{question}
Welke conclusie wordt rechtstreeks door deze informatie gedragen?

\begin{oplossing}
Regelmatig ontbijten en hogere schoolresultaten hangen in de onderzochte
groep samen.
\end{oplossing}
\end{question}

\begin{question}
Welke mogelijke derde variabelen kunnen een rol spelen?

\begin{oplossing}
Bijvoorbeeld thuissituatie, slaap, gezondheid, algemene leefgewoonten
of studieroutines.
\end{oplossing}
\end{question}

\begin{question}
Herschrijf de titel zodat hij overeenkomt met wat de gegevens werkelijk
aantonen.

\begin{oplossing}
Bijvoorbeeld:
\begin{quote}
``Leerlingen die regelmatig ontbijten behalen in dit onderzoek
gemiddeld hogere schoolresultaten.''
\end{quote}
\end{oplossing}
\end{question}

\end{oefening}




\begin{oefening}[Wiskunde en fysica: van correlatie naar verantwoorde claim]{\oefU \ESS}

Gebruik

\csv
{https://raw.githubusercontent.com/Ingmar-1986/wiskunde_cursus_2/main/modules/Data/Statistiek/wiskfys.csv}
{wiskfys.csv}

\begin{question}
Maak een spreidingsdiagram en bepaal $r$.

\begin{oplossing}
De puntenwolk is stijgend en ongeveer lineair. De correlatie is
\[
r\approx0{,}743.
\]
\end{oplossing}
\end{question}

\begin{question}
Geef twee plausibele factoren die zowel het resultaat voor wiskunde
als voor fysica kunnen beïnvloeden.

\begin{oplossing}
Bijvoorbeeld algemene studievaardigheid, voorkennis, motivatie,
studietijd of redeneervaardigheid. Andere goed gemotiveerde antwoorden
zijn mogelijk.
\end{oplossing}
\end{question}

\begin{question}
Formuleer een correcte eindconclusie die samenhang en causaliteit
duidelijk uit elkaar houdt.

\begin{oplossing}
In deze dataset bestaat een positieve lineaire samenhang tussen de
scores voor wiskunde en fysica. De gegevens tonen op zichzelf niet aan
dat een hogere score voor het ene vak de hogere score voor het andere
veroorzaakt.
\end{oplossing}
\end{question}

\end{oefening}



\begin{oefening}[Van observatie naar onderzoek]{\oefG}

Kies één van de datasets uit dit thema waarin een duidelijke samenhang
zichtbaar is.

Ontwerp een vervolgonderzoek waarmee je een mogelijke causale verklaring
sterker zou kunnen toetsen.

Beschrijf minstens:

\begin{enumerate}
  \item de causale hypothese;
  \item mogelijke derde variabelen;
  \item welke gegevens bijkomend nodig zijn;
  \item welke variabelen gecontroleerd of gerandomiseerd moeten worden;
  \item welke conclusie na het onderzoek nog steeds te sterk zou zijn.
\end{enumerate}

\begin{oplossing}
Er zijn meerdere correcte ontwerpen mogelijk. Een sterk antwoord maakt
duidelijk dat een correlatie aanleiding kan geven tot een hypothese,
maar dat een sterkere onderzoeksopzet nodig is om alternatieve
verklaringen uit te sluiten.
\end{oplossing}

\end{oefening}




\end{document}
